\Askhsh{31549_SOLUTION}{1}
\begin{ylh}{1}
    \item [Συναρτήσεις - Μελέτη Συνάρτησης]
\end{ylh}

\noindent\textbf{Λύση}

\begin{alist}
    \item [α)] Η $f$ είναι παραγωγίσιμη στο $(0,+\infty)$ με $f'(x) = \frac{1-\ln x}{x^2}$.
    \[ f'(x)=0 \iff \frac{1-\ln x}{x^2}=0 \iff 1-\ln x=0 \iff \ln x=1 \iff x=e \]
    \[ f'(x)>0 \iff \frac{1-\ln x}{x^2}>0 \iff 1-\ln x>0 \iff \ln x<1 \iff 0<x<e \]

    \begin{center}
    \begin{tabular}{|c|c||c|c|c|}
        \hline
        $x$ & $0$ & $e$ & $+\infty$ \\
        \hline\hline
        $f'(x)$ & $|$ & $+$ & $0$ & $-$ \\
        \hline
        $f(x)$ & $|$ & $\nearrow$ & $\circ$ & $\searrow$ \\
        \hline
    \end{tabular}
    \end{center}

    Συνεπώς η $f$ είναι γνησίως αύξουσα στο $(0,e]$, γνησίως φθίνουσα στο $[e,+\infty)$ και παρουσιάζει μέγιστο για $x=e$ το $f(e)= \frac{\ln e}{e} = \frac{1}{e}$.

    \item [β)] Επειδή $e<2022<2023$ και $f$ γνησίως φθίνουσα στο $[e,+\infty)$ είναι
    \[ f(2022)>f(2023) \iff \frac{\ln 2022}{2022} > \frac{\ln 2023}{2023} \]
    \[ 2023 \ln 2022 > 2022 \ln 2023 \]
    \[ \ln 2022^{2023} > \ln 2023^{2022} \]
    \[ 2022^{2023} > 2023^{2022} \]

    \item [γ)] Η $f'$ είναι παραγωγίσιμη στο $(0,+\infty)$ με
    \[ f''(x) = \frac{-2x-x(1-\ln x)}{x^4} = \frac{2\ln x -3}{x^3} \]
    \[ f''(x)=0 \iff \frac{2\ln x-3}{x^3}=0 \iff 2\ln x-3=0 \iff \ln x=\frac{3}{2} \iff x=e^{3/2} \iff x=\sqrt{e^3} \]
    \[ f''(x)>0 \iff \frac{2\ln x-3}{x^3}>0 \iff 2\ln x-3>0 \iff \ln x>\frac{3}{2} \iff x>e^{3/2} \iff x>\sqrt{e^3} \]

    \begin{center}
    \begin{tabular}{|c|c||c|c|c|}
        \hline
        $x$ & $0$ & $\sqrt{e^3}$ & $+\infty$ \\
        \hline\hline
        $f''(x)$ & $|$ & $-$ & $0$ & $+$ \\
        \hline
        $f(x)$ & $|$ & $\cap$ & $\oplus$ & $\cup$ \\
        \hline
    \end{tabular}
    \end{center}

    Συνεπώς η $f$ είναι κοίλη στο $(0, \sqrt{e^3}]$, κυρτή στο $[\sqrt{e^3}, +\infty)$ και παρουσιάζει καμπή για $x = \sqrt{e^3}$. Επειδή $f(\sqrt{e^3}) = \frac{\ln \sqrt{e^3}}{(\sqrt{e^3})^2} = \frac{3/2}{e^3} = \frac{3}{2e^3}$, το σημείο καμπής είναι το $(\sqrt{e^3}, \frac{3}{2e^3})$.

    \item [δ)] Από \text{Θ.Μ.Τ} για την $f$ στο $[2021,2022]$ υπάρχει $\xi_1 \in (2021,2022)$ με
    \[ f'(\xi_1) = \frac{f(2022)-f(2021)}{2022-2021} = f(2022)-f(2021) \]
    Από \text{Θ.Μ.Τ} για την $f$ στο $[2022,2023]$ υπάρχει $\xi_2 \in (2022,2023)$ με
    \[ f'(\xi_2) = \frac{f(2023)-f(2022)}{2023-2022} = f(2023)-f(2022) \]
    Επειδή $\sqrt{e^3} < 2021 < \xi_1 < 2022 < \xi_2$ και $f'$ γνησίως αύξουσα στο $[\sqrt{e^3}, +\infty)$ είναι
    \[ f'(\xi_1) < f'(\xi_2) \]
    \[ f(2022)-f(2021) < f(2023)-f(2022) \]
    \[ 2f(2022) < f(2023)+f(2021) \]
\end{alist}